\chapter*{Introdução}
\addcontentsline{toc}{chapter}{Introdução}
\label{cap:introducao}

% =====================================================================
% PAPEL DESTA INTRODUÇÃO:
%   Apresenta o problema, a pergunta, a hipótese, os objetivos, a
%   justificativa, o produto educacional e a estrutura do texto. NÃO
%   desenvolve conceitos: cada um deles tem casa própria nos capítulos
%   e aqui aparece apenas com a remissão correspondente.
% =====================================================================

Em poucos anos os modelos de linguagem de grande escala deixaram os laboratórios de pesquisa e
passaram a fazer parte da rotina de quem ensina. Ferramentas conversacionais como o ChatGPT, o
Gemini e o Claude estão a um clique de qualquer professor e respondem em segundos a pedidos que
antes custavam horas de trabalho, entre eles o de redigir questões para uma lista de exercícios
ou para uma avaliação. No ensino de Matemática essa promessa é especialmente atraente, porque a
elaboração de boas questões exige tempo, repertório e cuidado com detalhes que nem sempre cabem
na jornada de um professor da Educação Básica.

A mesma rapidez traz consigo um problema que esta dissertação toma como ponto de partida. Os
modelos de linguagem escrevem sobre Matemática com grande fluência e no entanto não garantem a
correção daquilo que escrevem. Erros de cálculo, gabaritos incompletos e resoluções que saltam
etapas aparecem redigidos com a mesma segurança das respostas corretas, de modo que nada no texto
avisa o leitor de que algo está errado. Esse fenômeno, que aqui se denomina \emph{erro
silencioso} e que é discutido na Seção~\ref{sec:erro-silencioso}, transfere ao professor uma
tarefa de revisão minuciosa que muitas vezes anula o ganho de tempo prometido pela ferramenta. Ao
lado do erro matemático há ainda um segundo risco menos visível, pois uma questão pode estar
correta e mesmo assim ser inadequada para o nível da turma ou não realizar a habilidade da Base
Nacional Comum Curricular que declara cobrir.

Diante desse cenário a pesquisa se organiza em torno de uma pergunta central. \textbf{Como apoiar
o professor de Matemática do Ensino Médio na geração de questões com modelos de linguagem
garantindo correção matemática verificável e adequação didática reconhecida?} A resposta proposta
parte da constatação de que nenhum componente isolado resolve o problema. O modelo de linguagem é
o único capaz de transformar uma especificação pedagógica em um enunciado plausível, mas é
justamente ele quem produz o erro silencioso. Por outro lado um sistema de computação algébrica
calcula com exatidão e não é capaz de redigir uma questão. Daí decorre a hipótese de trabalho,
argumentada por inteiro na Seção~\ref{sec:hipotese}, segundo a qual a combinação de um modelo de
linguagem gerador com um verificador simbólico determinístico, um agente crítico de natureza
didática e um orquestrador de ciclos de revisão produz questões mais confiáveis do que qualquer
desses componentes produziria sozinho.

Para examinar essa hipótese o trabalho tem como objetivo geral projetar, implementar e avaliar um
sistema multiagente para a geração assistida de questões de Matemática do Ensino Médio, com
mecanismos automáticos de verificação da correção matemática e de avaliação da qualidade
didática, submetendo as questões produzidas ao julgamento de professores em exercício. Esse
objetivo se desdobra nos seguintes objetivos específicos.

\begin{enumerate}
\item Caracterizar a partir da literatura as limitações dos modelos de linguagem em raciocínio
  matemático e as estratégias contemporâneas de mitigação.
\item Definir com base na BNCC do Ensino Médio e em referenciais didáticos consolidados uma
  especificação de entrada que permita ao professor descrever a questão que deseja.
\item Projetar uma arquitetura multiagente com papéis bem definidos de geração, verificação
  simbólica, crítica didática e orquestração.
\item Implementar um protótipo funcional em Python restrito a um recorte de habilidades do Ensino
  Médio.
\item Construir um banco curado de questões geradas e aprovadas pelo sistema e alinhadas a
  habilidades da BNCC.
\item Avaliar as questões geradas por meio de um painel de professores de Matemática.
\item Disponibilizar o sistema, o banco de questões e um manual de uso como produto educacional.
\end{enumerate}

O recorte adotado acompanha a unidade temática que a própria BNCC usa como exemplo de
reagrupamento de habilidades, a das funções polinomiais de primeiro e segundo graus, e a estende
por analogia às funções exponenciais, logarítmicas, trigonométricas e definidas por várias
sentenças e às progressões aritméticas e geométricas. Ao todo são quinze habilidades, apresentadas
na Seção~\ref{sub:bncc-habilidades}. A escolha não é apenas de conveniência computacional, já que
equações, funções e sequências admitem formalização simbólica e portanto conferência
determinística da resposta, ao passo que a unidade escolhida corresponde a um bloco de ensino que
o professor reconhece no próprio planejamento.

O desenvolvimento seguiu a abordagem da pesquisa em ciência do design, adequada ao mestrado
profissional por unir a construção de um artefato à produção de conhecimento sobre ele. O
trabalho foi conduzido em ciclos sucessivos de construção e avaliação, descritos no início do
Capítulo~\ref{cap:implementacao}, nos quais cada versão do sistema foi posta à prova com
provedores reais de modelos de linguagem antes de seguir para a etapa seguinte. Essa escolha
metodológica mostrou-se decisiva, uma vez que vários dos defeitos mais relevantes encontrados ao
longo da pesquisa só apareceram quando o sistema foi exercitado com gerações reais e não com os
testes automatizados.

A relevância do trabalho para o PROFMAT se manifesta em duas frentes. De um lado o produto
educacional resultante é diretamente aplicável à prática docente e foi concebido para ser usado
por professores que não programam. De outro o trabalho contribui para o debate sobre o uso
responsável da inteligência artificial na Educação Matemática, indo além da descrição das
ferramentas para examinar de que modo seus erros podem ser tornados visíveis ao professor. Em vez
de prometer um sistema que não erra, a dissertação propõe um sistema que declara com precisão o
que foi conferido em cada questão e o que continua dependendo do julgamento docente, posição que
se apoia na compreensão do professor como mediador indispensável discutida na
Seção~\ref{sec:papel-mediador}.

O produto educacional é composto por três entregáveis articulados, descritos no
Capítulo~\ref{cap:produto}. O primeiro é o próprio sistema multiagente, com uma interface web por
meio da qual o professor especifica, acompanha e avalia as questões geradas. O segundo é um banco
curado de questões aprovadas pelo ciclo completo do sistema, que pode ser usado mesmo por quem
não deseja instalar programa algum. O terceiro é o manual do professor, que orienta a instalação e
o uso e sobretudo explica em linguagem acessível os limites da ferramenta.

O texto está organizado em sete capítulos, além desta introdução e das considerações finais. O
Capítulo~\ref{cap:contexto} situa o uso da inteligência artificial na educação e no ensino de
Matemática e apresenta o erro silencioso do ponto de vista de quem ensina. O
Capítulo~\ref{cap:tecnica} reúne a fundamentação técnica sobre modelos de linguagem, sistemas
multiagentes e computação algébrica e termina com a formulação da hipótese arquitetural. O
Capítulo~\ref{cap:didatica} apresenta os referenciais didáticos que sustentam tanto a especificação
de entrada quanto os critérios de qualidade usados pelo sistema. Em seguida o
Capítulo~\ref{cap:arquitetura} descreve o projeto do sistema em nível conceitual e o
Capítulo~\ref{cap:implementacao} trata de sua implementação, incluindo custo, desempenho e
exemplos reais de execução. O Capítulo~\ref{cap:avaliacao} relata a avaliação das questões por
professores de Matemática da rede pública e o Capítulo~\ref{cap:produto} descreve o produto
educacional do ponto de vista de quem o recebe. Por fim as considerações finais retomam a pergunta
de pesquisa à luz dos resultados, discutem as limitações do estudo e apontam desdobramentos
possíveis.
